\begingroup
\centering
\begin{tikzpicture}[
  font=\small,
  box/.style={draw,rounded corners=2pt,align=center,inner sep=6pt,line width=.6pt},
  wide/.style={box,text width=12.5cm,fill=black!3},
  half/.style={box,text width=5.65cm,minimum height=2cm},
  flow/.style={-{Latex[length=2mm]},line width=.7pt},
  node distance=6mm and 7mm
]
\node[wide] (client) {\textbf{客户端：任务结束后构建主梯度子空间（PGS）}\\[3pt]
当前任务样本级梯度 $+$ 上一任务返回的历史子空间\\[2pt]
奇异值加权、拼接与截断分解 $\rightarrow$ 上传逐层基与奇异值};
\node[half,below=8mm of client.south west,anchor=north west,xshift=1mm] (uniform)
 {\textbf{服务器：FedSubMerge}\\[3pt]融合全部可用客户端方向\\加权基拼接与重新压缩\\返回共享保护子空间};
\node[half,right=of uniform] (adaptive)
 {\textbf{服务器：FedSubMerge-AD}\\[3pt]逐层计算子空间距离\\选择相关客户端并融合\\返回客户端特定保护子空间};
\draw[flow] (client.south -| uniform.north) -- (uniform.north);
\draw[flow] (client.south -| adaptive.north) -- (adaptive.north);
\node[wide,below=8mm of uniform.south east,anchor=north,xshift=3.5mm] (update)
 {\textbf{客户端：下一任务的投影训练}\\[3pt]
当前小批量梯度 $\rightarrow$ 保护子空间的正交补投影 $\rightarrow$ 优化器更新\\[2pt]
任务结束后，以当前梯度与返回的历史摘要递归更新 PGS};
\draw[flow] (uniform.south) -- (uniform.south |- update.north);
\draw[flow] (adaptive.south) -- (adaptive.south |- update.north);
\node[font=\footnotesize,align=center,text width=12.7cm,below=4mm of update]
 {常规轮次交换模型参数；任务结束后交换 PGS；不回放历史原始图像。};
\end{tikzpicture}
\par
\endgroup
